\section*{Chapter \ref{ch:b3:quantum-model-systems} --- Quantum Mechanics for Chemists: Model Systems}

\begin{solution}{exo:b3:quantum-model-systems:1}
$E_1 = h^2/8m_eL^2$. (a) $L = \qty{1.00}{nm}$: $E_1 = \qty{6.02e-20}{J} =
\qty{0.376}{eV}$; $\Delta E_{12} = 3E_1$ and $\lambda = hc/3E_1 =
\qty{1099}{nm}$, in the near infrared. (b) $L = \qty{0.50}{nm}$: $E_1$ is four
times larger, \qty{2.41e-19}{J} $= \qty{1.50}{eV}$, and $\lambda =
\qty{275}{nm}$, in the ultraviolet.
\end{solution}

\begin{solution}{exo:b3:quantum-model-systems:2}
$n_x^2 + n_y^2 + n_z^2 = 3$ (1,1,1): $g = 1$; $6$ (2,1,1 and permutations):
$g = 3$; $9$ (2,2,1): $g = 3$; $11$ (3,1,1): $g = 3$; $12$ (2,2,2): $g = 1$.
\end{solution}

\begin{solution}{exo:b3:quantum-model-systems:3}
$[\hat x^2,\hat p]f = -\iu\hbar x^2f' + \iu\hbar(x^2f)' = 2\iu\hbar xf$, so
$[\hat x^2,\hat p] = 2\iu\hbar\hat x$. $[\hat p^2,\hat x]f = -\hbar^2[(xf)'' -
xf''] = -2\hbar^2f'$, so $[\hat p^2,\hat x] = -2\iu\hbar\hat p$.
\end{solution}

\begin{solution}{exo:b3:quantum-model-systems:4}
$E_J/hc = BJ(J+1)$: 0, 21.19, 63.56 and \qty{127.12}{cm^{-1}}, with $g = 1$,
3, 5, 7. $I = h/(8\pi^2cB) = \qty{2.643e-47}{kg.m^2}$ (with $c$ in
\unit{cm/s}).
\end{solution}

\begin{solution}{exo:b3:quantum-model-systems:5}
$|\psi_n|^2$ is symmetric about $L/2$, so $\langle x\rangle = L/2$. With
$\sin^2u = \frac12(1 - \cos2u)$ and two integrations by parts,
$\langle x^2\rangle = L^2\big(\frac13 - \frac{1}{2n^2\pi^2}\big)$. For $n = 1$,
$\langle x^2\rangle = 0.283L^2$, less than the classical $L^2/3$: the ground
state piles up in the middle. As $n$ grows the classical value is recovered.
\end{solution}

\begin{solution}{exo:b3:quantum-model-systems:6}
$1\,\unit{cm^{-1}} = \qty{0.011963}{kJ/mol}$. ZPE(\ce{H2}) $= 2200.6 \times
0.011963 = \qty{26.33}{kJ/mol}$, ZPE(\ce{D2}) $= \qty{18.63}{kJ/mol}$; difference
\qty{7.69}{kJ/mol}. Same \omterm{def:b3:quantum-model-systems:oscillator}{force constant}, so $\tilde\omega \propto \mu^{-1/2}$
and the expected ratio is $\sqrt{\mu_D/\mu_H} = 1.4137$ (atomic masses
2.01410 and 1.00783); the measured ratio is 1.4127, within 0.07\,\% (the
small difference comes from the electrons, which do not follow the nuclei
perfectly).
\end{solution}

\begin{solution}{exo:b3:quantum-model-systems:7}
$\int_0^Lx^2(L - x)^2\,\dd x = L^5/30$, so $N = \sqrt{30/L^5}$. Then
\[
\langle\psi_1|\phi\rangle = \sqrt{\tfrac2L}\sqrt{\tfrac{30}{L^5}}\int_0^Lx(L -
x)\sin\tfrac{\pi x}{L}\,\dd x = \frac{\sqrt{60}}{L^3}\cdot\frac{4L^3}{\pi^3} =
\frac{4\sqrt{60}}{\pi^3} = 0.99928 .
\]
The parabola is 99.86\,\% ground state (the square of the overlap).
\end{solution}

\begin{solution}{exo:b3:quantum-model-systems:8}
$\langle f|\hat pg\rangle = -\iu\hbar\int f^*g'\,\dd x = -\iu\hbar[f^*g] +
\iu\hbar\int f^{*\prime}g\,\dd x = \int(-\iu\hbar f')^*g\,\dd x = \langle\hat
pf|g\rangle$, the bracket vanishing at the ends. For $\hat x\hat p$, using that $\hat x$ and $\hat p$ are each \omterm{def:b3:quantum-model-systems:hermitian}{Hermitian}:
$\langle f|\hat x\hat pg\rangle = \langle\hat xf|\hat pg\rangle = \langle\hat p\hat
xf|g\rangle$, and $\hat p\hat x = \hat x\hat p - \iu\hbar \ne \hat x\hat p$: not
\omterm{def:b3:quantum-model-systems:hermitian}{Hermitian} (its symmetrised form $\frac12(\hat x\hat p + \hat p\hat x)$ is).
\end{solution}

\begin{solution}{exo:b3:quantum-model-systems:9}
$L = 6 \times 139.7 = \qty{838}{pm}$, $N = 6$: $\lambda = 8m_ecL^2/(7h) =
\qty{331}{nm}$. The absorption is in the ultraviolet: hexatriene is
colourless.
\end{solution}

\begin{solution}{exo:b3:quantum-model-systems:10}
$\kappa = \sqrt{2m(V_0 - E)}/\hbar$ with $V_0 - E = \qty{0.20}{eV}$:
$\kappa_H = \qty{9.82e10}{m^{-1}}$, $\kappa_D = \qty{1.388e11}{m^{-1}}$. The
prefactors cancel in the ratio: $T_H/T_D = \eu^{2a(\kappa_D - \kappa_H)} =
\eu^{4.06} = 58$, the exact value: the barrier is thick for both
($\kappa a = 4.9$ and 6.9).
\end{solution}

\begin{solution}{exo:b3:quantum-model-systems:11}
Levels $E_m = m^2\hbar^2/2m_eR^2$: $m = 0$ holds two electrons, $m = \pm1$
four. The highest filled is $|m| = 1$, the lowest empty $|m| = 2$: $\Delta E =
3\hbar^2/2m_eR^2 = \qty{9.38e-19}{J}$, $\lambda = \qty{212}{nm}$, in the
ultraviolet, as benzene's strong absorption is.
\end{solution}

\begin{solution}{exo:b3:quantum-model-systems:12}
$\langle r\rangle = \frac{4\pi}{\pi a^3}\int_0^\infty r^3\eu^{-2r/a}\dd r =
\frac{4}{a^3}\cdot\frac{3!}{(2/a)^4} = \frac32a$, \qty{79.4}{pm} for $a = a_0$.
$\langle 1/r\rangle = \frac{4}{a^3}\cdot\frac{1!}{(2/a)^2} = \frac1a$, so
$\langle V\rangle = -e^2/(4\pi\varepsilon_0a)$. Since $E_1 =
-e^2/(8\pi\varepsilon_0a)$ (from $a = 4\pi\varepsilon_0\hbar^2/\mu e^2$),
$\langle V\rangle = 2E_1$, and the mean kinetic energy is $-E_1$.
\end{solution}

\begin{solution}{pb:b3:quantum-model-systems:1}
\textbf{1.} Each of the $j$ atoms gives one p orbital; the $j - 1$ carbons and
one nitrogen give one electron each and the other nitrogen its lone pair,
since the cation's charge is spread over the chain: $N = j + 1$, that is 6,
8 and 10.
\textbf{2.} Six electrons fill $n = 1$, 2, 3: HOMO $n = 3$, LUMO $n = 4$.
\textbf{3.} $\Delta E = E_{N/2+1} - E_{N/2} = (N+1)h^2/8mL^2$ and $\lambda =
hc/\Delta E = 8mcL^2/h(N+1)$.
\textbf{4.} $L_0 = 4l$, $6l$, $8l$: 559, 838 and \qty{1118}{pm}.
\textbf{5.} 147, 257 and \qty{374}{nm}.
\textbf{6.} All far too short (by factors 3.6, 2.3 and 1.9): the model's box is too
small, since $\lambda \propto L^2$; the electrons move over a longer
distance than the chain between the nitrogens.
\textbf{7.} With the origin at the centre, $\psi_n \propto \cos(n\pi u/L)$ for
odd $n$ and $\sin(n\pi u/L)$ for even $n$ ($u = x - L/2$): even and odd
functions of $u$.
\textbf{8.} If $n + n'$ is even, $\psi_n$ and $\psi_{n'}$ have the same parity;
their product times $u$ is odd and integrates to zero over a symmetric
interval.
\textbf{9.} HOMO $N/2$ and LUMO $N/2 + 1$ have $n + n' = N + 1$, odd: always
allowed.
\textbf{10.} Second dye: HOMO 4, LUMO 5. $4 \to 6$: sum even, forbidden.
$3 \to 5$: sum even, forbidden.
\textbf{11.} The next allowed one from $n = 4$ is $4 \to 7$, with
$\Delta E$ larger by $(49 - 16)/(25 - 16) = 3.67$: at $\lambda/3.67$,
\qty{164}{nm} for the second dye, deep in the ultraviolet.
\textbf{12.} It depends only on the symmetry of the potential about the
chain's centre, which the real symmetric dyes share.
\textbf{13.} $E = hc/\lambda = \qty{3.294e-19}{J} = \qty{2.06}{eV}$.
\textbf{14.} $hc\tilde\omega_e = \qty{0.371}{eV}$.
\textbf{15.} $2hcB = \qty{4.21e-22}{J} = \qty{2.63}{meV}$.
\textbf{16.} $kT = \qty{25.7}{meV}$.
\textbf{17.} Rotational levels only ($2hcB \ll kT$); vibrational ($15kT$) and
electronic ($80kT$) excitations are essentially not populated.
\textbf{18.} Electronic: visible and ultraviolet; vibrational: infrared;
rotational: microwaves (and far infrared).
\textbf{19.} $L = \sqrt{\lambda h(N+1)/8mc}$ with $N = 8$: $L = \qty{1283}{pm}$.
\textbf{20.} $\delta = (1283 - 838)/2 = \qty{222}{pm}$.
\textbf{21.} $L = 559 + 445 = \qty{1003}{pm}$, $\lambda = \qty{474}{nm}$, 9.5\,\%
below 524 nm.
\textbf{22.} $L = 1118 + 445 = \qty{1562}{pm}$, $\lambda = \qty{732}{nm}$, 2.9\,\%
above 711 nm.
\textbf{23.} Into the aromatic rings that carry the nitrogens: the $\pi$
system does not end at the nitrogen atoms, and the effective box includes
part of each ring.
\textbf{24.} $\delta/l = 222/139.7 = 1.6$: \textbf{the box extends about
\qty{222}{pm}, a bond and a half, beyond each nitrogen}, and with that one
length the model reproduces the three colours within 10\,\%.
\end{solution}
